\section{Name the bytes at each boundary}
\begin{figure}[htbp]
  \centering
  \begin{tikzpicture}[every node/.style={draw,rounded corners,align=center,font=\small},node distance=7mm,alt={A vertical pipeline joins PNG IDAT payloads, inflates their zlib stream, reconstructs filtered RGB rows, and packs framebuffer words.}]
    \node (file) {PNG file\\signature and chunks};
    \node[below=of file] (stream) {joined IDAT payloads\\one zlib stream};
    \node[below=of stream] (rows) {inflated bytes\\one filter byte plus samples per row};
    \node[below=of rows] (rgb) {reconstructed RGB bytes\\three bytes per pixel};
    \node[below=of rgb] (words) {display words\\one \texttt{0x00RRGGBB} word per pixel};
    \draw[-{Latex}] (file)--(stream);
    \draw[-{Latex}] (stream)--(rows);
    \draw[-{Latex}] (rows)--(rgb);
    \draw[-{Latex}] (rgb)--(words);
  \end{tikzpicture}
  \caption{Four transformations with inspectable outputs.}
\end{figure}
For the first image, each row has six sample bytes. The inflated representation has
seven bytes per row because it also stores a filter selector. Five rows therefore
occupy 35 inflated bytes, while the reconstructed samples occupy 30 bytes. Packing
the ten pixels uses 40 bytes if the destination is tightly packed. A display with a
512-word stride still advances 2048 bytes between row starts. These are four different
sizes; keep their names in the program rather than calling all of them the image size.
